\documentclass[12pt,twoside]{article}
\usepackage{fancyhdr,makeidx,times,multicol,geometry}
\usepackage{keyval,ifthen,mparhack,manyfoot,poemscol}
\newcounter{scounter}
\newcommand\step{\stepcounter{scounter}\Roman{scounter}}

\title{Postal del Cielo}
\author{J.J. Gómez Cadenas}
\begin{document}
\maketitle

\poemtitle{Primera Ciudad}
\begin{poem}
\begin{stanza}
La primera ciudad nace mucho antes,\\
de que pises París por vez primera,\\
hecha está de materia de los sueños.
\end{stanza}
\begin{stanza}
Y en tu mente, avenidas extendiéndose\\
ríos de mármol y plata, arrogantes\\
e imperiales palacios contemplándote.
\end{stanza}
\begin{stanza}
Tú, soñador, que las calles recorres,\\
tu imaginación llena de palabras,\\
leídas en romances y novelas\\
repetidas por poetas y cantores,\\
filmadas en escenas recurrentes.
\end{stanza}
\begin{stanza}
Imaginada, la ciudad primera,\\
que visitaste una vez, pero en vano,\\
tratarás de volver, sin encontrarla.
\end{stanza}
\begin{stanza}
Como Ítaca, París no existió,\\
fue conjurada por deseo y niebla,\\
por delirios imperiales y polvo,\\
por su sucio río, sus catacumbas,\\
sus piedras vacías y tu nostalgia.
\end{stanza}
\end{poem}

\poemtitle{Segunda Ciudad}
\begin{poem}
\begin{stanza}
Tampoco visitaste la ciudad segunda,\\
aquella por la que vagabas a los veinte,\\
enamorado hasta los huesos, sin saber,\\
que París y tu amor los inventabas.
\end{stanza}
\begin{stanza}
Los largos días, paseando,\\
por calles y plazas recién estrenadas,\\
las noches febriles,\\
en las que casi tocabais el Nirvana.\\
Los cafés que os hacían sentiros artistas,\\
el sándwich robado, jugando a la bohemia,\\
y la luna brillando,\\
sobre vuestros interminables besos.
\end{stanza}
\begin{stanza}
Sin comprender que todo era,\\
solo el guion de una vulgar opereta,\\
la ciudad invadida por gente que repetía,\\
sin saberlo, las mismas líneas,\\
paseando por los mismos parques,\\
y besándose, como un disciplinado ejército\\
bajo cada puente del Sena.
\end{stanza}
\begin{stanza}
Nada de eso veías,\\
cegado por la luz de tu escasa veintena,\\
nunca, a esa edad, a París viajaste,\\
pudo ser Viena, o Roma, Tel Aviv o Praga,\\
solo fuiste a la ciudad de tu juventud.
\end{stanza}
\end{poem}

\poemtitle{La Ciudad real}
\begin{poem}
\begin{stanza}
Hay un día en que vuelves con tus hijos,\\
a caminar por las calles de antaño,\\
y todos los arcángeles del cielo,\\
parecen envolverte con sus alas.
\end{stanza}
\begin{stanza}
Recorres París de nuevo con ellos,\\
tu niña de once años lleva el mapa,\\
que os guiará en el subterráneo hostil\\
y así llegaréis a Disney a salvo.
\end{stanza}
\begin{stanza}
Ella se ocupa. Mamá le ha explicado,\\
que papá está medio ciego y no sabe,\\
navegar por los océanos del metro.\\
Así que ase tu mano firmemente.
\end{stanza}
\begin{stanza}
Pero aún más firme es la mano de tu hijo\\
de siete años. Mamá le ha encargado,\\
protegerte de todos los peligros.\\
De acero su mirada, aterroriza,\\
a los duendes y trolls que os acechan.
\end{stanza}
\begin{stanza}
Y están de nuevo tus ojos abiertos,\\
y el corazón del pecho se te sale,\\
y es que este amor en nada se parece,\\
a los fuegos fatuos que antes seguías.
\end{stanza}
\begin{stanza}
Este amor tan hondo y que tanto arde,\\
con sus llamas a la ciudad esculpe,\\
haciéndola real por vez primera.
\end{stanza}
\end{poem}
\end{document}
